% !PRINTED 311
\weberpage{311}
%<<CONT
the beginning of the expansion for $w_z$:
\weq{4}{w_z = q^{-\frac{2}{5}} e^{\frac{c}{2}\frac{\pi i}{5}} + \cdots}
From this it follows that the power sums of the $w_z$, up to the fourth inclusive, are constants, since they cannot contain even the first power of $j(\omega)$, and that the product of the $w_z$ must be a \gesperrt{linear} function of $j(w)$, % [sic] German source prints w here, but omega four words later.
in which $j(\omega)$ has the coefficient 1.

The equation of the $w$ therefore has the form:
\weq{5}{w^5 + b_1 w^4 + b_2 w^3 + b_3 w^2 + b_4 w + b_5 = j(\omega),}
where the $b$ are rational numbers. These coefficients can be calculated by carrying the expansion (4) further. Here we take another route, which, strictly speaking, belongs to the complex multiplication treated in the next part, but which, being so simple, may nevertheless quite well find its place here.

If we set $\omega = i = \sqrt{-1}$, then
\[ \omega = -\frac{1}{\omega} \]
and consequently [§ 34, (11), (12), (14)]:
\[ f_1(i) = f_2(i) = \sqrt[8]{2}, \quad f(i) = \sqrt[4]{2}, \]
hence [§ 54, (5)]:
\weq{6}{\gamma_3(i) = 0, \quad \gamma_2(i) = 12, \quad j(i) = 12^3.}
But one finds further, from the transformation formulae for $\eta(\omega)$ [§ 34, (4), (5)]:
\[ \eta(\omega+1) = e^{\frac{\pi i}{12}} \eta(\omega), \quad \eta\left(-\frac{1}{\omega}\right) = \sqrt{-i\omega}\,\eta(\omega), \]
if the decomposition
\[ 5 = (2+i)(2-i) \]
is taken into account:
\[ \eta\left(\frac{12+i}{5}\right) = e^{\frac{\pi i}{6}} \eta\left(\frac{2+i}{5}\right) = e^{\frac{\pi i}{6}} \eta\left(\frac{1}{2-i}\right) = \sqrt{1+2i}\,\eta(i), \]
hence
\[ v_2 = 1+2i, \quad v_3 = 1-2i, \]
and now, since for $\gamma_2 = 12$ two roots of equation (1) are known, one can easily decompose the left part of this equation into factors for this particular value of $\gamma_2$:
\[ v^6 + 10v^3 - 12v + 5 = (v^2+2v+5)(v^2+v-1)^2, \]
%CONT>>
